% !TeX root = part5.tex
\documentclass[../final.tex]{subfiles}
\begin{document}

% Условия: Семинар_5.pdf, 02.10.26.
% Задача 1: расшифровка двух страниц Seminar_5_by_iScanner.pdf.
% В рукописи заголовок «Семинар 4», но условие соответствует семинару 5.
% Задача 2: самостоятельное дополнение по присланному условию.
\practice{5}{02.10.2026}{Бозе-газ во внешнем поле}

\rtheory{Конденсация идеального бозе-газа}
Как и ранее, полагаем $k_B=1$: температура измеряется в единицах
энергии. Рассматриваем невзаимодействующие бозоны без внутреннего
вырождения. Энергию и химический потенциал отсчитываем от основного
уровня. Тогда $\mu\leq0$, а средние числа заполнения равны
\[
 \overline n(\varepsilon)=\frac{1}{e^{(\varepsilon-\mu)/T}-1}.
\]
Отдельно выделяем основное состояние: $N=N_0+N_{\mathrm{ex}}$.
В квазиклассическом пределе сумму по возбуждённым уровням заменяем
интегралом с плотностью состояний $D(\varepsilon)$.
Введём обозначения
\[
 z=e^{\mu/T},\qquad
 g_s(z)=\sum_{l=1}^{\infty}\frac{z^l}{l^s},\qquad
 g_s(1)=\zeta(s),\qquad
 \frac{d g_s(z)}{dz}=\frac{g_{s-1}(z)}{z}.
\]
Используем интеграл Бозе
\[
 \int_0^\infty\frac{x^{s-1}\,dx}{z^{-1}e^x-1}
 =\Gamma(s)g_s(z).
\]
При $T<T_0$ в термодинамическом пределе $\mu=0$: возбуждённые
уровни вмещают лишь часть частиц, остальные находятся в конденсате.
Энергия основного состояния после выбранного сдвига равна нулю.

\task[1]{Бозе-газ в гармонической ловушке}
\condition
Бозе-газ из $N$ частиц находится в поле
$U(r)=m\omega^2r^2/2$. Найти температуру конденсации $T_0$,
теплоёмкость при $T<T_0$ и её скачок в точке конденсации.

\subsection*{Плотность состояний и температура конденсации}
Уровни трёхмерного осциллятора до сдвига нуля энергии имеют вид
\[
 E_n=\hbar\omega\left(n+\frac32\right),\qquad
 d_n=\frac{(n+1)(n+2)}2,\qquad n=0,1,\ldots
\]
При $T\gg\hbar\omega$ существенны большие $n$, поэтому
$\varepsilon=\hbar\omega n$ и
\[
 D(\varepsilon)=\frac{\varepsilon^2}{2(\hbar\omega)^3}.
\]
Следовательно,
\[
 N_{\mathrm{ex}}=\frac{1}{2(\hbar\omega)^3}
 \int_0^\infty\frac{\varepsilon^2\,d\varepsilon}
 {z^{-1}e^{\varepsilon/T}-1}
 =\left(\frac{T}{\hbar\omega}\right)^3g_3(z).
\]
В точке перехода $z\to1$, $N_{\mathrm{ex}}\to N$. Отсюда
\[
 \boxed{T_0=\hbar\omega\left(\frac{N}{\zeta(3)}\right)^{1/3}},
 \qquad \frac{N_0}{N}=1-\left(\frac{T}{T_0}\right)^3
 \quad(T<T_0).
\]

\subsection*{Энергия и теплоёмкость ниже перехода}
Учитывается полная энергия движения в ловушке, включая потенциал:
\[
 E=\int_0^\infty\frac{\varepsilon D(\varepsilon)\,d\varepsilon}
 {z^{-1}e^{\varepsilon/T}-1}
 =\frac{3T^4}{(\hbar\omega)^3}g_4(z).
\]
При $T<T_0$ имеем $z=1$, поэтому
\[
 \boxed{C_V\equiv\left(\frac{\partial E}{\partial T}\right)_{N,\omega}
 =12\zeta(4)\left(\frac{T}{\hbar\omega}\right)^3
 =12N\frac{\zeta(4)}{\zeta(3)}\left(\frac{T}{T_0}\right)^3}.
\]

\subsection*{Скачок теплоёмкости}
Выше перехода конденсата нет. Дифференцируя условие постоянства $N$,
получаем
\[
 0=\frac{3T^2}{(\hbar\omega)^3}g_3(z)
 +\frac{T^3}{(\hbar\omega)^3}\frac{g_2(z)}{z}\frac{dz}{dT},
 \qquad \frac1z\frac{dz}{dT}=-\frac3T\frac{g_3(z)}{g_2(z)}.
\]
Тогда
\[
 C_V(T>T_0)=N\left[12\frac{g_4(z)}{g_3(z)}
 -9\frac{g_3(z)}{g_2(z)}\right].
\]
В частности, поскольку $\mu=T\ln z$,
$\left.d\mu/dT\right|_{T_0+}=-3\zeta(3)/\zeta(2)$,
как и в рукописи. Односторонние пределы равны
\begin{align*}
 C_V(T_0-)&=12N\frac{\zeta(4)}{\zeta(3)}\simeq10{,}805N,\\
 C_V(T_0+)&=N\left[12\frac{\zeta(4)}{\zeta(3)}
 -9\frac{\zeta(3)}{\zeta(2)}\right]\simeq4{,}228N.
\end{align*}
\answer
Определяя скачок при нагревании как значение выше перехода минус
значение ниже, находим
\[
 \boxed{\Delta C_V=C_V(T_0+)-C_V(T_0-)
 =-9N\frac{\zeta(3)}{\zeta(2)}\simeq-6{,}577N}.
\]
При конечном $N$ переход сглажен; эти формулы относятся к ведущему
квазиклассическому приближению большого числа частиц.

\clearpage
\task[2]{Бозе-газ в линейном поле над непроницаемой плоскостью}
\condition
Потенциал равен $U(z)=Fz$ при $z>0$ и $U(z)=\infty$ при $z<0$,
где $F>0$. Найти $T_0$ и теплоёмкость при $T<T_0$.

\subsection*{Геометрия и плотность состояний}
Решение этой задачи добавлено самостоятельно. Вдоль плоскости газ
занимает большую конечную площадь $S$; эффекты её краёв не учитываем.
Ответ выражается через поверхностную плотность $N/S$. Без фиксации
$S$ или $N/S$ численное значение температуры определить нельзя.
В полуклассическом приближении
\begin{align*}
 D(\varepsilon)
 &=\frac{S}{(2\pi\hbar)^3}\int_0^\infty dz\int d^3p\,
 \delta\left(\varepsilon-\frac{p^2}{2m}-Fz\right)\\
 &=\frac{S(2m)^{3/2}}{4\pi^2\hbar^3}
 \int_0^{\varepsilon/F}(\varepsilon-Fz)^{1/2}\,dz\\
 &=\frac{S(2m)^{3/2}}{6\pi^2\hbar^3F}\varepsilon^{3/2}.
\end{align*}
Отсюда, используя $\Gamma(5/2)=3\sqrt\pi/4$, получаем
\[
 N_{\mathrm{ex}}=A T^{5/2}g_{5/2}(z),\qquad
 A=\frac{S}{F}\left(\frac{m}{2\pi\hbar^2}\right)^{3/2}.
\]
Здесь $z=e^{\mu/T}$ обозначает активность, а в потенциале $z$
является пространственной координатой.
Ту же формулу даёт интегрирование локальной концентрации:
\[
 N_{\mathrm{ex}}=\frac{S}{\lambda_T^3}
 \int_0^\infty g_{3/2}(e^{(\mu-Fz)/T})\,dz
 =\frac{ST}{F\lambda_T^3}g_{5/2}(e^{\mu/T}),\quad
 \lambda_T=\sqrt{\frac{2\pi\hbar^2}{mT}}.
\]

\subsection*{Конденсация и теплоёмкость}
При $\mu=0$ число возбуждённых частиц конечно. Из условия
$N=A T_0^{5/2}\zeta(5/2)$ находим
\[
 \boxed{T_0=\left[\frac{NF}{S\zeta(5/2)}
 \left(\frac{2\pi\hbar^2}{m}\right)^{3/2}\right]^{2/5}},
 \qquad \frac{N_0}{N}=1-\left(\frac{T}{T_0}\right)^{5/2}.
\]
Поскольку $\Gamma(7/2)/\Gamma(5/2)=5/2$, полная энергия равна
\[
 E=\frac52 A T^{7/2}g_{7/2}(e^{\mu/T}).
\]
Это сумма кинетической энергии и энергии в поле:
$E_{\mathrm{kin}}=(3/2)AT^{7/2}g_{7/2}$,
$E_{\mathrm{pot}}=AT^{7/2}g_{7/2}$.
Ниже перехода $\mu=0$, и при фиксированных $N,S,F$
\[
 \boxed{C_V=\frac{35}{4}A\zeta(7/2)T^{5/2}
 =\frac{35}{4}N\frac{\zeta(7/2)}{\zeta(5/2)}
 \left(\frac{T}{T_0}\right)^{5/2}}.
\]

\subsection*{Область применимости и единицы}
Для линейного потенциала характерный масштаб квантования по высоте
$\varepsilon_F^{\mathrm{lin}}=(\hbar^2F^2/(2m))^{1/3}$.
Полученный ответ справедлив при
$\varepsilon_F^{\mathrm{lin}}\ll T\lesssim T_0$ и малых интервалах
уровней поперечного движения по сравнению с $T$.
Это стандартное трёхмерное квазиклассическое приближение.
Если устремлять только $S\to\infty$ при фиксированных $F,N/S$,
точный нижний уровень движения по высоте образует двумерную ветвь:
в строгом таком пределе истинной конденсации идеального газа при
конечной температуре нет. Поэтому формулу $T_0$ следует понимать
в указанном приближении для большой конечной системы либо в
пределе ослабления ловушки, сохраняющем трёхмерное описание.

Для температуры в кельвинах в формулах нужно заменить $T$ на $k_BT$:
температуры $T_0$ в ответах делятся на $k_B$, а теплоёмкости
умножаются на $k_B$ при записи через отношение $T/T_0$.

\end{document}
